\section{Dualidad y topologías débiles}

\begin{ejercicio}
Sea $X$ un espacio de Banach, $\{x_n\} \subset X$, $\{f_n\} \subset X^*$ tales que $\{x_n\}$ converge débilmente a $x \in X$ y $\{f_n\}$ converge en norma a $f \in X^*$. Prueba que $\lim\limits_n f_n(x_n) = f(x)$.
\end{ejercicio}

\begin{ejercicio}
En $\ell_\infty$ consideremos la siguiente sucesión de funcionales: para $n \in \bb{N}$, $\varphi_n \in \ell_\infty^*$ está definido por
\[
\varphi_n(x) = x(n) \quad \forall x \in \ell_\infty
\]
Prueba que $\{\varphi_n : n \in \bb{N}\} \subset B_{X^*}$ pero ninguna subsucesión de la sucesión $\{\varphi_n\}$ converge en la topología débil-$*$, $\sigma(\ell_\infty^*, \ell_\infty)$. ¿Contradice esto el Teorema de Banach-Alaoglu?
\end{ejercicio}

\begin{ejercicio}
Consideremos los funcionales $\varphi_n \in C\left[-1, 1\right]^*$ dados por
\[
\varphi_n(f) = f\left(\nicefrac{-1}{n}\right) - f\left(\nicefrac{1}{n}\right) \quad \forall f \in C\left[-1, 1\right], \; \forall n \in \bb{N}
\]
Prueba que $\{\varphi_n\} \longrightarrow 0$ en la topología $\sigma(C\left[-1, 1\right]^*, C\left[-1, 1\right])$ pero no en la norma de $C\left[-1, 1\right]^*$.
\end{ejercicio}

\begin{ejercicio}
Sea, para cada $n \in \bb{N}$, $\varphi_n: c_0 \longrightarrow \bb{K}$ definida por
\[
\varphi_n(x) = \frac{x(1) + x(2) + \dots + x(n)}{n} \quad \forall x \in c_0
\]
\begin{enumerate}[label=\textbf{\alph*)}]
    \item Prueba que $\varphi_n \in c_0^*$ y calcula su norma.
    \item Prueba que $\{\varphi_n\} \longrightarrow 0$ en la topología $\sigma(c_0^*, c_0)$ pero no en la topología $\sigma(c_0^*, c_0^{**})$.
\end{enumerate}
\end{ejercicio}

\begin{ejercicio}
Sea $X$ un espacio de Banach separable de dimensión infinita. Construye un operador lineal y continuo $T: X \longrightarrow c_0$ que no alcance su norma. [Idea: prueba que existe una sucesión $\{x_n^*\} \subset S_{X^*}$ que converge a $0$ en la topología $\sigma(X^*, X)$].
\end{ejercicio}